\section*{第 \ref{ch:g1:plants-around-us} 章——身边的植物}

\begin{solution}{exo:g1:plants-around-us:1}
根（在土壤下）、茎、叶、花。
\end{solution}

\begin{solution}{exo:g1:plants-around-us:2}
(a) 根。(b) 叶。(c) 花。(d) 茎。
\end{solution}

\begin{solution}{exo:g1:plants-around-us:3}
它由一粒种子长成，它会取食（从土壤中吸水，从天空获得光），
它会生长，会结出种子，这些种子长成新的雏菊，
它终有一天会死亡。生命的各种特征它都具备，走路
则不在其中。
\end{solution}

\begin{solution}{exo:g1:plants-around-us:4}
树的茎是树干；\omterm{def:g1:plants-around-us:parts}{叶子}
着生于\omterm{def:g1:plants-around-us:tree}{分支}上。
\end{solution}

\begin{solution}{exo:g1:plants-around-us:5}
水、光，以及供根生长的地方——通常是土壤。
\end{solution}

\begin{solution}{exo:g1:plants-around-us:6}
是浇水。\omterm{def:g1:plants-around-us:plant}{植物}因缺水而萎蔫；
浇水几小时后，\omterm{def:g1:plants-around-us:parts}{叶子}会重新挺立。
\end{solution}

\begin{solution}{exo:g1:plants-around-us:7}
相同之处：两者都是\omterm{def:g1:plants-around-us:plant}{植物}，都有根、茎、叶和花，
都需要水和光。不同之处：树的茎是木质的树干，
而雏菊的茎柔软、呈绿色，树
要大得多，寿命也长得多。
\end{solution}

\begin{solution}{exo:g1:plants-around-us:8}
窗边的\omterm{def:g1:plants-around-us:plant}{植物}应是绿色、挺立的；柜子里的\omterm{def:g1:plants-around-us:plant}{植物}
苍白、细长、虚弱，尽管已经浇过水。缺少的
是光——只有水是不够的。
\end{solution}

\begin{solution}{exo:g1:plants-around-us:9}
例如：\omterm{def:g1:plants-around-us:plant}{植物}在生长——每周都比前一周高——
并且会转向或倾向光源，像向日葵跟随太阳一样。
生长慢并不等于什么都没做。
\end{solution}

\begin{solution}{exo:g1:plants-around-us:10}
是根。“地下”指的是图中根所在的位置——
我们吃的胡萝卜是一根又大又肥的橙色根，
\omterm{def:g1:plants-around-us:plant}{植物}把食物储存在其中。
\end{solution}
